\section{Conclusion}
\label{sec:conclusion}
This paper presented a reproducible diagnostic protocol for sparse-concept and calibration evaluation in Knowledge Tracing, combining train-only KC-frequency stratification with an explicit strict cold-start group, learner-based and temporal splits, per-stratum calibration analysis (ECE and the Brier decomposition, with reliability diagrams), sample-size-aware reliability flags, and a reproducible audit workflow (L1--L8) that operationalizes standard safeguards plus a predictive-sanity check. The protocol is intended as reusable evaluation infrastructure for KT research, in particular for future graph-based and self-supervised approaches whose claimed gains concentrate precisely on sparse KCs, where evaluation is least stable.

The instantiation on three datasets illustrates what the protocol surfaces that aggregate metrics do not. Rather than a universal AUC failure, we observe substantial predictive heterogeneity. The strongest empirical finding is a dataset-dependent calibration susceptibility: while population-level metrics remained competitive, calibration often degraded severely as KC training evidence decreased in specific datasets. For example, on ASSISTments 2012, SimpleKT's ECE increased consistently from $0.113$ on dense concepts to $0.158$ on medium concepts and $0.225$ on sparse concepts (subject to a Limited reliability flag, $N = 403$). A controlled downsampling of training rows for the same originally dense KCs produced a within-KC ECE increase only in some dataset--model cells (large for SimpleKT on Junyi Academy; absent for DKT and SimpleKT on ASSISTments 2012; Section~\ref{sec:a9_sparsify}), so observational frequency--calibration gradients should not be read as a universal frequency dose--response. Sparse-concept diagnostics are most informative when a dataset contains a meaningful low-frequency tail, those strata have sufficient evaluation support, and calibration reliability varies systematically with training evidence or deployment involves concept cold-start (Section~\ref{sec:a8_when_needed}). Occupancy reporting remains useful even when that empirical priority is low. A standardized, dataset-agnostic protocol is therefore a reporting framework whose \emph{claim intensity} is conditional, not a requirement that every KT log be narrated as a sparse-calibration failure. Finally, channel L8 provided the empirical demonstration that the workflow is more than a preventative log: it caught a prediction--label alignment artifact that would otherwise have been misread as genuine cold-start failure, and a controlled fault-injection validation later measured sensitivity to additional alignment-fault classes (Appendix~\ref{app:l8_fault}). L8 is reported as sensitive to the classes we tested, not as a detector of every evaluation bug. L1--L7 remain standard safeguards made inspectable; they are not offered as a new auditing methodology.

Future work can extend this diagnostic foundation by evaluating across multiple temporal seeds or cutoffs to capture variance, by incorporating a classical baseline that exploits item difficulty for unseen learners under learner-based splits, and by applying the protocol to diagnose graph-augmented and self-supervised KT architectures.
